% Wird aus 05_Programmkonzept.tex eingebunden
\subsection{Belegung der Taster}
\label{sec:bedienkonzept}
Zur Bedienung gibt es nur die zwei Taster des AC2398. 
Damit kann der Bediener den Zeiger korrigieren und das Programm nach einem Fehler neu starten. 
Der rote Taster ist ein Öffner, sein Signal wird deshalb vor dem Baustein umgedreht.

\begin{table}[H]
\centering\small
\begin{tabular}{|L{3.6cm}|L{10.4cm}|}
\hline
\textbf{Bedienung} & \textbf{Funktion} \\ \hline
Grün länger als 300\,ms halten & Zeiger fährt nach oben (positiver Offset), bis der Taster losgelassen oder das Skalenende erreicht wird. Grüne LED leuchtet. \\ \hline
Rot länger als 300\,ms halten & Zeiger fährt nach unten (negativer Offset), bis der Taster losgelassen oder das Skalenende erreicht wird. Rote LED leuchtet. \\ \hline
Beide Taster 3\,s halten & Neustart mit Referenzfahrt: Offset wird gelöscht, der Zeiger fährt auf 0 und anschließend auf die aktuelle Temperatur. Beide LEDs blinken mit 2\,Hz, solange die Referenzfahrt läuft. Ein anliegender Fehler wird dabei quittiert. \\ \hline
\end{tabular}
\caption{Belegung der Taster des AC2398}
\end{table}

\subsection{Bedeutung der LEDs}
Die LEDs zeigen über schnelles oder langsames Blinken an, was das Programm gerade macht. 
Wird ein Taster gedrückt, leuchtet seine LED dauerhaft. Das gilt nicht während des Starts, der Referenzfahrt und im Fehlerzustand.

\begin{table}[H]
\centering\small
\begin{tabular}{|L{2.6cm}|L{2.6cm}|L{8.8cm}|}
\hline
\textbf{LED grün} & \textbf{LED rot} & \textbf{Bedeutung} \\ \hline
blinkt 2\,Hz & blinkt 2\,Hz & Referenzfahrt läuft \\ \hline
blinkt 1\,Hz & aus & Normalbetrieb, Anzeige aktuell \\ \hline
blinkt 1\,Hz & blinkt 1\,Hz & Temperatur außerhalb 0--36\,\textdegree C, Zeiger am Skalenende \\ \hline
aus & blinkt 1\,Hz & Fehler, Motor gestoppt \\ \hline
\end{tabular}
\caption{Bedeutung der LED-Zustände}
\end{table}

Eine kompakte Anleitung für den Bediener findet sich in Kapitel~\ref{sec:bedienungsanleitung}.
